\section{Introduction}

Gaussian quadrature \index{Gaussian quadrature} is a numerical integration method that achieves exact results for polynomials up to degree $2n-1$ using only $n$ function evaluations. Gauss developed this method in his work on the arithmetic-geometric mean and the theory of special functions.

In this chapter, we cover:
\begin{itemize}
    \item Orthogonal polynomials and their zeros
    \item Gauss-Legendre \index{Gauss-Legendre quadrature} quadrature
    \item Gauss-Hermite \index{Gauss-Hermite quadrature}, Gauss-Laguerre \index{Gauss-Laguerre quadrature}, Gauss-Chebyshev \index{Gauss-Chebyshev quadrature} quadrature
    \item The Golub-Welsch algorithm \index{Golub-Welsch algorithm} for computing nodes and weights
    \item Error analysis and convergence
    \item Multidimensional quadrature and sparse grids
\end{itemize}

\section{Orthogonal Polynomials}

Let $w(x) > 0$ be a weight function on $[a,b]$. The monic orthogonal polynomials \index{orthogonal polynomial} $p_n(x)$ satisfy
\[
\int_a^b p_n(x) p_m(x) w(x) dx = 0 \quad (n \ne m).
\]
They obey a three-term recurrence \index{three-term recurrence}:
\[
p_{n+1}(x) = (x - \alpha_n) p_n(x) - \beta_n p_{n-1}(x).
\]

\section{Gaussian Quadrature}

\begin{theorem}[Gauss Quadrature]
Let $x_1, \ldots, x_n$ be the zeros of $p_n(x)$. There exist positive weights $w_1, \ldots, w_n$ such that
\[
\int_a^b f(x) w(x) dx = \sum_{i=1}^n w_i f(x_i)
\]
for all polynomials $f$ of degree $\le 2n-1$.
\end{theorem}

The weights are $w_i = \int_a^b \frac{p_n(x)}{(x-x_i)p_n'(x_i)} w(x) dx$.

\section{Classical Gauss Quadrature Rules}

\begin{table}[H]
\centering
\begin{tabular}{llll}
\toprule
Rule & Interval & Weight & Polynomials \\
\midrule
Gauss-Legendre & $[-1,1]$ & $1$ & Legendre $P_n$ \\
Gauss-Hermite & $(-\infty,\infty)$ & $e^{-x^2}$ & Hermite $H_n$ \\
Gauss-Laguerre & $[0,\infty)$ & $e^{-x}$ & Laguerre $L_n$ \\
Gauss-Chebyshev 1st & $[-1,1]$ & $(1-x^2)^{-1/2}$ & Chebyshev $T_n$ \\
Gauss-Chebyshev 2nd & $[-1,1]$ & $(1-x^2)^{1/2}$ & Chebyshev $U_n$ \\
\bottomrule
\end{tabular}
\end{table}

\section{Golub-Welsch Algorithm}

Compute nodes and weights via the symmetric tridiagonal matrix
\[
J = \begin{pmatrix}
\alpha_0 & \sqrt{\beta_1} & & \\
\sqrt{\beta_1} & \alpha_1 & \sqrt{\beta_2} & \\
& \ddots & \ddots & \ddots \\
& & \sqrt{\beta_{n-1}} & \alpha_{n-1}
\end{pmatrix}
\]
The eigenvalues of $J$ are the nodes $x_i$. The weights are $w_i = \mu_0 v_{i,1}^2$ where $v_i$ is the normalized eigenvector and $\mu_0 = \int w(x) dx$.

\section{Python Implementation}

In \texttt{python/gauss\_quadrature.py}:

\begin{lstlisting}[style=python, caption={Key functions in python/gauss\_quadrature.py}]
gauss_legendre(n)           # Nodes, weights for Legendre
gauss_hermite(n)            # Hermite
gauss_laguerre(n)           # Laguerre
gauss_chebyshev(n, kind)    # Chebyshev 1 or 2
golub_welsch(alpha, beta)   # General Golub-Welsch
integrate(f, a, b, n)       # Gauss-Legendre quadrature
gauss_kronrod(n)            # Kronrod extension \index{Kronrod quadrature} for error estimate
\end{lstlisting}

\section{Numerical Examples}

\begin{example}[Gauss-Legendre]
\begin{lstlisting}[style=python]
>>> from gauss_quadrature import gauss_legendre, integrate
>>> x, w = gauss_legendre(5)
>>> print("Nodes:", x)
>>> print("Weights:", w)
>>> # Integrate x^8 on [-1,1] (exact for degree <= 9)
>>> val = integrate(lambda x: x**8, -1, 1, 5)
>>> print(val, "exact:", 2/9)
\end{lstlisting}
\end{example}

\section{Exercises}

\begin{exercise}
Derive the recurrence coefficients $\alpha_n, \beta_n$ for Legendre polynomials.
\end{exercise}

\begin{exercise}
Implement Gauss-Radau and Gauss-Lobatto quadrature (include endpoints).
\end{exercise}

\begin{exercise}
Use Gauss-Hermite quadrature to compute $\int_{-\infty}^{\infty} e^{-x^2} f(x) dx$ for various $f$.
\end{exercise}

\begin{exercise}
Compare Gaussian quadrature with Clenshaw-Curtis for analytic functions.
\end{exercise}



\section{Exercise Solutions}

\begin{solution}
Gauss-Legendre quadrature with $n$ nodes integrates polynomials of degree $2n-1$ exactly. The nodes are the roots of the Legendre polynomial\index{Legendre polynomial} $P_n(x)$, and the weights are $w_i = \frac{2}{(1-x_i^2)[P_n'(x_i)]^2}$. For $\int_{-1}^{1} f(x)\,dx \approx \sum_{i=1}^n w_i f(x_i)$, the error is $O(f^{(2n)}(\xi))$.
\end{solution>

\begin{solution}
The Golub-Welsch algorithm computes Gauss quadrature nodes and weights by finding the eigenvalues and first components of eigenvectors of the symmetric tridiagonal Jacobi matrix. For weight function $w(x) = 1$ on $[-1,1]$, the Jacobi matrix has entries $a_k = 0$, $b_k = \frac{k}{\sqrt{4k^2-1}}$.
\end{solution>

\begin{solution}
Composite Gauss quadrature divides $[a,b]$ into $m$ subintervals and applies $n$-point Gauss-Legendre on each. The error is $O(m \cdot h^{2n}) = O((b-a)h^{2n-1})$ where $h = (b-a)/m$. For smooth integrands, this achieves spectral accuracy with few nodes per interval.
\end{solution>

\begin{solution}
Gauss-Kronrod quadrature adds $n+1$ optimally chosen nodes to an $n$-point Gauss rule, giving a higher-order estimate for error control. The Clenshaw-Curtis rule uses Chebyshev polynomial\index{Chebyshev polynomial} nodes and is stable for adaptive quadrature. Both are implemented in numerical libraries like SciPy's \texttt{quad}.
\end{solution>


\section{References}

\begin{itemize}
    \item \cite{gauss-works} -- Gauss's \textit{Methodus nova integralium valores per approximationem inveniendi}
    \item \cite{golub-welsch} -- Golub \& Welsch (1969), \textit{Calculation of Gauss Quadrature Rules}
    \item \cite{stroud-secrest} -- Stroud \& Secrest, \textit{Gaussian Quadrature Formulas}
    \item \cite{trefethen-approx} -- Trefethen, \textit{Approximation Theory and Approximation Practice}
\end{itemize}